\documentclass[10pt,a4paper]{amsart}
\usepackage[margin=19mm]{geometry}
\usepackage{amsmath,amssymb,mathtools}
\usepackage[scr=rsfs,cal=euler]{mathalfa}
\usepackage{xfrac}
\usepackage[unicode,hidelinks,bookmarks=false]{hyperref}
\newcommand{\ie}{i.e.\ }
\newcommand{\cf}{cf.\ }
\newcommand{\resp}{respectively\ }
\newcommand{\Subsection}{\subsection}
\providecommand{\nobreakdash}{\nobreak\hbox{-}}
\setlength{\emergencystretch}{2em}
\multlinegap0pt
\usepackage{tikz}
\usetikzlibrary{cd}
\usepackage{xcolor}
\usepackage{amssymb}
\usepackage{mathtools}
\usepackage[shortlabels]{enumitem}
\usepackage{tikz}
\newtheorem{non}{Non}
\newtheorem{prop}{Proposition}
\newcommand{\A}{\mathbf{A}}
\newcommand{\Bl}{\mathbf{Bl}}
\newcommand{\Oh}{\mathcal{O}}
\newcommand{\ProS}{\mathbf{P}}
\newcommand{\QCoh}{\mathcal{D}}
\newcommand{\Sm}{\text{Sm}}
\newcommand{\Smoo}{\text{Sm}}
\newcommand{\Unit}{\textbf{1}}
\newcommand{\V}{\mathbf{V}}
\newcommand{\from}{\colon}
\newcommand{\iso}{\simeq}
\title{A six-functor formalism for syntomic cohomology}
\author{Niklas Kipp}
\date{Source: arXiv:2608.15744v1 (CC BY 4.0)}
\begin{document}
\maketitle

\noindent\emph{Source and licence.} A six-functor formalism for syntomic cohomology. Version arXiv:2608.15744v1 (CC BY 4.0), distributed by arXiv under the Creative Commons Attribution 4.0 International licence, http://creativecommons.org/licenses/by/4.0/, as recorded in the arXiv licence metadata for this exact version and shown on the abstract page https://arxiv.org/abs/2608.15744v1. Full licence notice: \texttt{licenses/arxiv-2608.15744-CC-BY-4.0.txt}.\par\smallskip

\noindent\textit{Editorial scope.} Counted unit: the article's prop prop (source0.tex lines 1101--1114). The article's complete original proof of it is delivered with it (source0.tex lines 1116--1147, one \texttt{proof} environment of 32 lines). No mathematics is written, repaired or re-derived: the statement and the proof are the article's own lines, in the article's own notation, and no retained line is substituted or re-flowed.  Retained uncounted as the source-local context the proof needs: lines 1054--1069.  Nothing was omitted from any retained span, so the package carries no omission record.  No retained line contains a citation, so the wrapper carries no bibliography and no bibliography entry.  The retained lines contain the article's own diagram code, which the wrapper typesets with the standard tikz packages; no image file is delivered and the package declares no asset dependency.  Editorial formatting only, in addition: an AMS \texttt{amsart} preamble that restates the article's own declarations and macros used by the retained lines, namely the environments \texttt{non}, \texttt{prop} on separate counters, together with the standard packages those lines need.  The retained environments print in this document as Non 1, Proposition 1 in order of appearance, whereas the article prints the same items with the numbering of its own full document.  The complete native source of the article as distributed in the arXiv e-print for this exact version is bundled unchanged as \texttt{sources/207777/source0.tex}.
\begin{non}\label{blowup squares}
   For $S\in \Smoo$ we write $\Sm^{\text{geom}}_{S}$ for the category of geometrically
   smooth morphisms over $S$. Furthermore, we will call a blow-up square 
      % https://q.uiver.app/#q=WzAsNCxbMSwxLCJYIl0sWzAsMSwiWiJdLFsxLDAsIlxcQmxfe1h9KFopIl0sWzAsMCwiRSJdLFsyLDBdLFsxLDBdLFszLDFdLFszLDJdXQ==
\[\begin{tikzcd}
	E & {\Bl_{Z}(X)} \\
	Z & X
	\arrow[from=1-1, to=1-2]
	\arrow[from=1-1, to=2-1]
	\arrow["p", from=1-2, to=2-2]
	\arrow["i"', from=2-1, to=2-2]
	\arrow["f", from=1-1, to=2-2]
\end{tikzcd}\]
   associated to a regular closed immersion $i\from Z\to X$ between objects in $\Sm^{\text{geom}}_{S}$ a \emph{smooth 
   blowup square}. 
\end{non}

\begin{prop}\label{blowup excision prop}
	Consider an object $S\in \Sm$. Then any blow-up square in $\Sm_{S}^{\text{geom}}$
	as in \ref{blowup squares} induces a (co)cartesian square 
	    	   % https://q.uiver.app/#q=WzAsNCxbMCwwLCJcXFVuaXRfe1h9Il0sWzEsMCwiaV97XFxhc3R9aV57XFxhc3R9XFxVbml0X3tYfSJdLFswLDEsInBfe1xcYXN0fXBee1xcYXN0fVxcVW5pdF97WH0iXSxbMSwxLCJmX3tcXGFzdH1mXntcXGFzdH1cXFVuaXRfe1h9Il0sWzAsMV0sWzAsMl0sWzIsM10sWzEsM11d
\[\begin{tikzcd}
	{\Unit_{X}} & {i_{\ast}i^{\ast}\Unit_{X}} \\
	{p_{\ast}p^{\ast}\Unit_{X}} & {f_{\ast}f^{\ast}\Unit_{X}}
	\arrow[from=1-1, to=1-2]
	\arrow[from=1-1, to=2-1]
	\arrow[from=1-2, to=2-2]
	\arrow[from=2-1, to=2-2]
\end{tikzcd}\]
   in $\QCoh(X)$.
\end{prop}

\begin{proof}
	By pulling back along the structure map (using proper base change), we can deduce the 
	assertion for blow-up squares of the form 
	  	       % https://q.uiver.app/#q=WzAsNCxbMSwxLCJYIl0sWzAsMSwiWiJdLFsxLDAsIlxcQmxfe1h9KFopIl0sWzAsMCwiRSJdLFsyLDBdLFsxLDBdLFszLDFdLFszLDJdXQ==
\[\begin{tikzcd}
	\ProS_{Z}^{n-1} & \V_{\ProS^{n-1}_{Z}}(\Oh(1)) \\
	Z & \A^{n}_{Z}
	\arrow[from=1-1, to=1-2]
	\arrow[from=1-1, to=2-1]
	\arrow["p", from=1-2, to=2-2]
	\arrow["0"', from=2-1, to=2-2]
	\arrow["f", from=1-1, to=2-2]
\end{tikzcd}\]
   for general $Z\in \Sm$. For a general smooth blow-up square, we can use assertion $(f)$ in the 
   definition of geometrically smooth morphisms to Zariski locally on $X$ find cartesian squares 
      % https://q.uiver.app/#q=WzAsNixbMCwwLCJaIl0sWzAsMSwiWCJdLFsxLDEsIlUiXSxbMSwwLCJaIl0sWzIsMSwiXFxBXntufV97Wn0iXSxbMiwwLCJaIl0sWzUsNCwiMCJdLFsyLDQsImUiLDJdLFszLDVdLFszLDJdLFsyLDEsImoiXSxbMywwXSxbMCwxLCJpIiwyXV0=
\[\begin{tikzcd}
	Z & Z & Z \\
	X & U & {\A^{n}_{Z}}
	\arrow["i"', from=1-1, to=2-1]
	\arrow[from=1-2, to=1-1]
	\arrow[from=1-2, to=1-3]
	\arrow[from=1-2, to=2-2]
	\arrow["0", from=1-3, to=2-3]
	\arrow["j", from=2-2, to=2-1]
	\arrow["e"', from=2-2, to=2-3].
\end{tikzcd}\]
   Let us write $Q_{X}, Q_{U}$ and $Q_{\A^{n}_{Z}}$ for the cohomology squares induced by the corresponding blow-up squares. 
   Then by open descent, to check that $Q_{X}$ is cartesian, it suffices to see that 
      $$j^{\ast}(Q_{X})\iso Q_{U} \iso e^{\ast}(Q_{\A^{n}_{Z}})$$
	is cartesian. But this we did see in the first half of the argument. 
\end{proof}

\end{document}
